\section{The probability of a dynamical cutoff}
\label{sec:cutoff}

The Gaussian limit proved in \Cref{prop:cumulant-clt} concerns the
pasted dynamics. To pass to the true dynamics, we couple both systems
using the same initial configuration. Mere convergence of the
discrepancy probability to zero would transfer uncentered observations,
but need not control their exact expectations on the scale
\(\mu^{-1/2}=\eps\). The stronger bound below supplies this control
in \Cref{lem:center-transfer}. It concerns the unobserved particle
system; all dynamical parameters are chosen independently of an
observation list or a moment order.

\begin{proposition}[Accuracy of the pasted dynamics]\label{prop:cutoff}
There is a sufficiently large dimensional integer \(\Gamma_0\) with the
following property.  Fix \(\Gamma\geq\Gamma_0\).  There is
\(L_0=L_0(\Gamma,f_0,f,T,\beta)\) such that, for any fixed \(L\geq L_0\),
the true dynamics and the pasted dynamics of
\Cref{thm:marked-expansion}, coupled using the same initial configuration,
satisfy
\begin{equation}\label{eq:cutoff-probability}
 \P\{\text{the two trajectories differ at some time in }[0,T]\}
 \leq C\eps^{1+\eta}
\end{equation}
for all sufficiently small \(\eps\), where one may take
\(\eta=1/20\).  Here the size cutoffs are the \(A_i,\Lambda_i\) of
\Cref{thm:dhm-package}, and \(C\) and the upper bound on \(\eps\) may
depend on the displayed fixed parameters.
\end{proposition}

We distinguish the geometric cycle excess \(r(M)\) of a graph from its
multislab complexity \(\rho(M)\).  A graph with \(b\) binary-contact
vertices, \(n\) particle lines, and \(c\) connected components has
\begin{equation}\label{eq:cutoff-euler}
 r(M)=b-n+c
\end{equation}
when each line has at least one vertex.  Equivalently, regarding contacts
as vertices and intervening particle segments as bonds,
\(r(M)=\#\{\text{bonds}\}-b+c\).
The graph's directed order is the order generated by increasing time on
each particle line.  Orders on disjoint lines need not be comparable.
All graph integrals below have the activity normalization of
\Cref{thm:marked-expansion}.

\subsection{Terminal witnesses and the reduction to bounded complexity}

Suppose the two systems first disagree at an incoming contact in slab
\(i\).  Immediately before this contact the systems agree.  If two
collision components are about to merge, each has size at most
\(\Lambda_i\) and excess at most \(\Gamma\).  Their union, including the
new contact, has size at most \(2\Lambda_i\) and excess at most
\(2\Gamma\).  If the contact is internal, the size does not change and
the new excess is at most \(\Gamma+1\leq2\Gamma\).  Thus the first
disagreement has a witness consisting of a connected, realizable
hard-sphere collision history \(M_{\mathrm{top}}\), stopped at this
contact, satisfying
\begin{equation}\label{eq:cutoff-witness}
 n(M_{\mathrm{top}})\leq2\Lambda_i,\qquad
 r(M_{\mathrm{top}})\leq2\Gamma,\qquad
 n(M_{\mathrm{top}})>\Lambda_i
 \ \text{or}\ 
 r(M_{\mathrm{top}})>\Gamma .
\end{equation}
After the stopping contact its lines may simply be extended without
further prescribed contacts.  In particular the stopping time need not
be a slab endpoint.

For every initial configuration the indicator of a first disagreement
is bounded by the number of these witnesses.  For a specified
\(n\)-particle witness, the factorial integral at the previous interface
has coefficient \(\mu^n/n!\); it has no prescribed top root and no
additional factor \(\mu^{-1}\).  Removing conditions involving particles
outside the witness only increases this positive bound.
Apply \Cref{cor:terminal-witness}.  This expands the preceding pasted
factorial measure with its exact activity and factorial coefficients.
Every lower component is attached to the terminal history through
particle lines.  The internal realizability restriction on
\(M_{\mathrm{top}}\) remains imposed.  Remainders are smaller than any
fixed power of \(\eps\).

Choose a fixed complexity threshold \(K\), after \(L\) has been fixed.
By \Cref{cor:terminal-witness,thm:dhm-package}, the sum of the positive
majorants with \(\rho>K\), including the activity cost of an artificial
top root when required, is \(O(\eps^3)\) if \(K\) is large enough.
The terms left to estimate therefore have
\begin{equation}\label{eq:cutoff-low-complexity}
 \rho\leq K .
\end{equation}
Their numbers of interface crossings, initial-link vertices, and
within-slab cycles are bounded.  Their total geometric cycle excess is
bounded as well.  We may discard initial-link restrictions in these
positive majorants; the remaining trajectory attachments to the
terminal component persist.

It remains to bound these low-complexity witnesses. We first localize
their birth energy, then find enough local geometric gains in the
terminal collision component, and finally sum its ordinary contacts
together with the attached lower histories.

\subsection{Energy localization}

The improved local estimates will be used only on bounded velocity
regions. We first remove very high birth energy by a crude forward
integration, without using those estimates.

\begin{lemma}[Removal of very high birth energy]\label{lem:cutoff-energy-tail}
For the positive graph majorants in
\eqref{eq:cutoff-low-complexity}, there is a fixed power \(H_*\)
such that restricting the total birth energy
\(P=\sum_{j\text{ a birth}}|v_j|^2\) to
\[
 P\leq P_{\eps}:=(1+|\log\eps|)^{H_*}
\]
changes their sum by \(O(\eps^M)\) for every fixed \(M\).
\end{lemma}

\begin{proof}
All atom and birth counts are bounded by a fixed power
\(B_\eps=(1+|\log\eps|)^H\), by the layer cutoffs and
\eqref{eq:cutoff-low-complexity}.
A crude count of all graphs, time orders, and coefficient
choices in this finite range is at most
\(\exp(C_*B_\eps^2\log(2+B_\eps))\).
For example, label all births, list at most \(B_\eps\) ordered
contacts by their two labels and C/O type, and then list their
slabs and interface gluings.  Discarding factorial denominators
and allowing all such lists gives this upper bound.  Bounded
complexity initial coefficients cost at most an additional
fixed power of \(B_\eps\).

The two-incoming-flight identity requires no velocity
localization.  Indeed, for fixed relative intercept \(x\)
and velocity \(v\), a transversal root of
\(|x+tv|=\eps\) has derivative
\(v\cdot(x+tv)/\eps\).  Thus
\[
 \int \delta(|x+tv|-\eps)
       [-v\cdot(x+tv)/\eps]_+\,\dd t\leq1:
\]
a straight flight has at most one incoming crossing of
a ball, and the flux cancels the scalar coarea Jacobian.
The outgoing-position and outgoing-velocity delta factors
then make a deterministic substitution with unit mass.
Zero-speed and grazing cases carry no transversal incoming
flux; use the same almost-everywhere convention, or exhaust
the strictly transversal set.  This proves the unrestricted
incoming-contact bound directly, before localization.

For a fixed fully ordered graph, let \(h_j\) be the absolute
input at birth \(j\): it is either \(|g_{i_j}|\) or an \(f_0\)
factor from an initial cumulant bound. Both have the common
Gaussian envelope from \Cref{prop:one-root,ass:kinetic}.
Integrate the original coarea measure forward, starting with
all birth phase points.
At each contact its two incoming affine lines have already
been determined.  The unrestricted identity just proved
integrates the contact time and outgoing
variables to an indicator at most one and a deterministic
update.  Births at later interfaces are inserted at their
assigned times; they can equivalently be integrated in advance
and held inactive until then.  Thus, after dropping all
encounter restrictions and initial links, a graph with \(n\)
births has mass at most
\[
 \mu^n\int\prod_{j=1}^n h_j(x_j,v_j)
       \ind_{\{\sum_j|v_j|^2>P_\eps\}}
       \prod_j\dd x_j\,\dd v_j
 \leq \mu^n C^n e^{-cP_\eps}.
\]
The last inequality uses half of the uniform birth Gaussian
for the energy tail and integrates the other half.  This
forward bound is an estimate of the original positive
measure, so it involves no local recollision gain.
Summing the crude finite count, and using \(n\leq B_\eps\),
gives an upper bound
\[
 \exp\{C_*B_\eps^2\log(2+B_\eps)
            +2B_\eps|\log\eps|-cP_\eps\}.
\]
Choose \(H_*>2H+2\).  The result is smaller than every
fixed power of \(\eps\).
\end{proof}

\subsection{Conditional local kernels and admissible cuts}

A \emph{cut} exposes a set of contact variables. Its shared lines
become fixed phase-space data for the unexposed graph: each specifies
a velocity and an affine free-flight position. A later piece is
integrated conditional on all variables exposed earlier in the cut
order. The available degree of an atom is its number of nonfixed edges,
counting internal bonds at each endpoint and counting free ends.

A history is regular if at every O-atom each fixed end is serial
with a free end; that serial
fixed/free pair is called a simple pair. A full history is regular.
A one-atom piece of type \(\{2\}\) has two nonfixed ends; its two fixed
ends are on the same time side if it is C, and are nonserial if it is O.
Types \(\{3\}\) and \(\{4\}\) have respectively one and zero fixed
ends. A \(\{33\}\) piece consists of two atoms joined by one bond,
each of degree three. It is type A if one atom can be cut as a type
\(\{3\}\) singleton and the other thereby becomes an allowed
\(\{2\}\) piece; it is type B
otherwise. A \(\{44\}\) piece has two degree-four atoms joined by one
bond. These conventions are
\cite[Definitions 8.2 and 8.11]{DHM2025}.

Cutting a pure-C set as free makes each crossing bond a free end of the
cut piece and a fixed end of the complement. At O-atoms, the serial
intercepts are equal, so the operation must preserve that equality.
Take each maximal serial segment whose internal atoms are O. Break its
bonds, reconnect its atoms belonging to the cut set in their order,
and supply free ends at its first and last cut atoms whenever the
segment continues outside the set. On the complement side, the extreme
remaining incidence receives the corresponding fixed end, while an
intermediate O-atom receives a simple pair. Every duplicated fixed
value equals the original intercept on that segment. This is the
regular cutting operation of \cite[Definition 8.3]{DHM2025}.

The associated kernels compose exactly: if \(A\) is cut first as free,
then
\begin{equation}\label{eq:cut-composition}
 I_M=I_A\circ I_{M\setminus A},
\end{equation}
with the repeated boundary variables identified as above. Thus the
variables of an earlier cut piece are exterior variables and are held
fixed while the later kernel is estimated. This includes any external
time or position used in a later support restriction. Equality
\eqref{eq:cut-composition} is
\cite[Proposition 8.14]{DHM2025}; it also follows by integrating the
serial delta functions in \eqref{eq:atom-O}. Splitting into finitely
many indicator cases is permitted. Deleting an O-atom means integrating
out its incoming encounter count and dropping the resulting indicator
bounded by one, as in \cite[Proposition 8.13]{DHM2025}.

\begin{lemma}[Conditional elementary bounds]\label{lem:dhm-local}
Consider the positive kernel \eqref{eq:edge-integral} of a regular
piece, with its fixed ends prescribed. Suppose the support has
\(|v_e|\le |\log\eps|^{C_*}\), and each \(x_e\) lies in a prescribed
ball of radius \(|\log\eps|^{C_*}\). The centers of those balls and
all exterior variables may be arbitrary. Then:
\begin{enumerate}[label=(\roman*)]
 \item A \(\{2\}\) piece has kernel mass at most one. For a
 \(\{3\}\) piece the kernel is parameterized by
 \begin{equation}\label{eq:local-three-parameters}
   \dd t\,\dd\omega\,\dd w\,
         |(v-w)\cdot\omega|,
 \end{equation}
 with \((t,\omega,w)\in\R\times S^2\times\R^3\) and the
 appropriate incoming hemisphere. Here \(v\) is the fixed-end
 velocity. If that end is outgoing, use the inverse elastic
 parametrization. All other states are determined by these
 parameters and the fixed end. A \(\{4\}\) piece has in addition
 the measure \(\mu\,\dd x\,\dd v\).
 Consequently their masses are respectively bounded by
 \(|\log\eps|^{C_*}\) and
 \(\mu|\log\eps|^{C_*}\). A neutral type-A \(\{33\}\) bound is
 obtained by a \(\{3\}\) integration followed by a \(\{2\}\)
 integration.
 \item A \(\{3\}\) piece gains \(\eps^\upsilon\) over its neutral
 bound if its support satisfies any one of
 \begin{equation}\label{eq:good-three-conditions}
 \begin{split}
 |t-t_*|&\le\eps^\upsilon,\qquad
 |v_e-v_{e'}|\le\eps^\upsilon,\\
 |x_e-x_*|&\le\eps^{1-\upsilon},\qquad
 |x_e-x_{e'}|\le\eps^{1-\upsilon}.
 \end{split}
 \end{equation}
 The two-edge conditions require distinct edges and, for O-atoms,
 exclude serial pairs. In the external-position condition \(e\)
 cannot be the fixed edge, or, for an O-atom, its serial free edge.
 The quantities \(t_*,x_*\) are fixed exterior parameters. The bound
 is uniform in their values.
 \item For a type-B \(\{33\}\) piece, both atoms are C. Suppose
 \(|t_1-t_2|\ge\eps^\upsilon\), and at each of its atoms every pair
 of distinct incident edges satisfies
 \(|v_e-v_{e'}|\ge\eps^\upsilon\). Its kernel mass is at most
 \(\eps^\upsilon|\log\eps|^{C_*}\), uniformly in its two fixed
 ends.
\end{enumerate}
In all three assertions, multiplying by an arbitrary bounded function
on the stated support costs its supremum. There is no assumption that
a fixed-end velocity has an independent Gaussian distribution.
\end{lemma}

\begin{proof}
Part (i) is the incoming coarea formula
\cite[Proposition 9.1, Equations (9.1)--(9.6)]{DHM2025}; the two-end
kernel counts at most one incoming crossing between the two fixed
flights. The powers of \(\mu\) also follow directly from
\eqref{eq:edge-integral}. The type-A assertion follows from
\eqref{eq:cut-composition} and its defining cut. Part (ii) is
\cite[Proposition 9.2(1), Equations (9.14)--(9.15)]{DHM2025}, with
precisely the restrictions on distinct and serial edges written there.
Part (iii) is \cite[Proposition 9.3(2), Equation (9.23)]{DHM2025};
its range \(0<\upsilon<1/(6d)\) contains \(1/81\) for \(d=3\).
The estimates are bounds for a positive kernel with a supremum over
all fixed ends, which gives the last assertion and permits their
successive conditional use.
\end{proof}

\begin{remark}\label{rem:conditional-extraction}
One may first move the ordinary \(\dd t\,\dd\omega\,\dd w\)
variables in \eqref{eq:local-three-parameters} outside later kernels
and then apply \Cref{lem:dhm-local} with them fixed. This is an
application of Tonelli's theorem to the positive coarea integrals and
\eqref{eq:cut-composition}; all support restrictions must remain until
the projected ordinary-parameter domain has been estimated. The source
uses exactly this order in \cite[Equations (9.66) and (9.74)]{DHM2025}.
The local estimates do not control the volume of that projected
domain. We will bound it using the common birth energy and causal
time order in \Cref{lem:cutoff-ellipsoid,lem:cutoff-time}, after proving
the extraction statement in \Cref{lem:cutoff-extraction}.
\end{remark}

A \emph{uniformly good} singleton or type-B pair will mean one of the
restricted kernels in \Cref{lem:dhm-local} with conditional mass
\(\eps^\upsilon|\log\eps|^{C_*}\).  All its gain-producing restrictions
are part of that kernel; they are not discarded later.
In particular, the phrase uniformly good includes the full velocity
and time nondegeneracy hypotheses required for a type-B pair.  It
does not mean merely that a pair has two distinct velocities.

The following is the precise operational part of the cutting
algorithms that we import.  It contains no improved probability
estimate.

\begin{proposition}[Pure-collision cutting alternatives, imported]\label{prop:cutoff-import}
Let \(M\) be a connected, full, realizable graph containing only
scattering contacts.  Full means that all its end variables are
integrated.  Define a transversal set \(A\subset M\) to be a set for
which \(M\setminus A=A^-\sqcup A^+\), with no directed bond from
\(A^+\) to \(A\cup A^-\), or from \(A\) to \(A^-\).
Let \(X_0(A)\) consist of vertices outside \(A\) with two bonds to
\(A\), and let \(X(A)\) be the closure obtained by repeatedly adding
vertices with two bonds to the set already grown.  The following
statements hold; in parts \textup{(ii)--(iv)}, \(A\) is assumed to be
connected.
\begin{enumerate}[label=\textup{(\roman*)}]
\item\label{item:cutoff-growth}
There is a chain of connected transversal sets
\(A_1,\ldots,A_m=M\), with \(A_1\) a tree, such that
\(X_0(A_j)\neq\varnothing\) for \(j<m\),
\[
 A_j\cup X(A_j)\subset A_{j+1},\qquad
 r(A_{j+1})=r(A_j\cup X(A_j)),\qquad
 r(A_{j+1})-r(A_j)\leq10|X(A_j)|.
\]
The intermediate growth from \(A_j\cup X(A_j)\) to \(A_{j+1}\)
adds one bond at each new vertex.
For any \(D\) transitions in this chain there is a cutting sequence
of all of \(M\) with one type \(\{4\}\), no type \(\{44\}\) or
type-B pair, and at least \(D\) type-A pairs.
\item\label{item:cutoff-bfk}
One can choose \(A^\pm\) such that each component of \(A^+\) is
attached to \(A\).  The two parts
\(X^\pm(A)=X(A)\cap A^\pm\) are equivalently obtained by adding
vertices with two predecessors, or two successors, in the grown set.
They are convex along particle lines.  If \(m_\pm\) is the number
of bonds from \(A\) to \(X^\pm(A)\), then
\[
 |X^\pm(A)|\leq G(m_\pm),
 \qquad
 G(q)=(2q+1)\binom q2(32q^{3/2})^{q^2}\quad(q\geq2).
\]
The harmless values \(G(0),G(1)\) may be replaced by one.
\item\label{item:cutoff-trans}
Put \(z=|X_0^\pm(A)|\) and \(a=r(A)\).  There is a cutting
sequence of \(M\) having one type \(\{4\}\), no type \(\{44\}\)
or type-B pair, and at least \(z/2-a-1\) type-A pairs.
\item\label{item:cutoff-main}
For this same choice of side there is a finite Borel partition
of the contact-parameter domain, with a subcase-dependent integer
\(y\geq0\), and the following two available sequences in each
subcase:
\begin{itemize}
\item a sequence with one type \(\{4\}\), no type \(\{44\}\)
or type-B pair, and at least \(y/10-3(a+1)\) uniformly good
singletons;
\item a sequence with no type \(\{44\}\), at most
\(S=a+z+y\) type \(\{4\}\) singletons, all type-B pairs
uniformly good, and at least
\[
 \frac{m_\pm-10^5S}{10}
\]
type-A or type-B pairs in total.
\end{itemize}
\end{enumerate}
The singleton and pair types not counted for a gain have the neutral
kernels described above.  On graphs of polylogarithmic size and
bounded excess the indicated partitions and operation sequences
have at most \(C^b|\log\eps|^{C_*}\) subcases.  The exponential
base \(C\) depends on dimensional choices, not on \(L\).
The required local restrictions at each uniformly good piece involve
its own variables and data exposed earlier in the cutting order.
\end{proposition}

\begin{proof}[Source identification]
The transversal definition and closure properties are
\cite[Definition 15.3 and Proposition 15.4]{DHM2025}.
The chain and its transition cutting sequence are the construction
in the proof of \cite[Proposition 15.1, case (1)]{DHM2025}.
The displayed bound \(G\), including its realizability hypothesis,
is \cite[Proposition 15.2]{DHM2025}; for \(X^\pm(A)\) its
entering boundary consists of the stated connection bonds.
That proposition uses the hard-sphere collision bound of
Burago, Ferleger and Kononenko~\cite{BuragoFerlegerKononenko1998},
adapted to extended dynamics in \cite[Appendix A]{DHM2025}.
Part \ref{item:cutoff-trans} is
\cite[Proposition 15.6]{DHM2025}.
The two sequences in \ref{item:cutoff-main} are
\cite[Proposition 15.10, (15.7)--(15.8)]{DHM2025}.
We deliberately weaken the first count in (15.7) when retaining
only good singletons.  Its proof has at most \(r(A)\) type-two
or type-A pieces inside \(A\), covering at most \(2r(A)\)
vertices, and one free seed.  The remaining weakly degenerate
vertices or later-cut members of weak pairs give good
singletons.  Discarding these exceptional vertices and weakening
the count gives the stated \(y/10-3(r(A)+1)\).
For the second sequence, the type-B pairs lie in the remaining
part of \(A\), after the weakly degenerate atoms and weak pairs
have been removed. Both the time separation and all the
velocity separations required by \Cref{lem:dhm-local}(iii)
therefore hold there.
The integer \(y\) is the size of the set \(Y_1\) in that
proposition's parameter partition.  Thus its name and internal
classification are unnecessary for the operational statement here.
The uniform local estimates and conditioning convention are
\Cref{lem:dhm-local}.  We use none of the final exponents in
DHM's truncation-error proposition.
\end{proof}

We also record a simple neutral completion argument; it will be useful
both after a seed cut and below a terminal component.

\begin{lemma}[Neutral completion]\label{lem:cutoff-neutral}
Suppose an exposed set is order-convex in the directed contact graph.
Every remaining component touching this set can be cut using only
type \(\{2\}\) and type \(\{3\}\) singletons.
The same conclusion holds with given top ends in place of the
exposed set.  Scattering and nonscattering contacts are both allowed.
\end{lemma}

\begin{proof}
If a remaining vertex is a successor of the exposed set, choose one
minimal among the unexposed successors.  It is adjacent to the exposed
set, and adding it preserves order convexity.  Every fixed incident
edge is on its incoming side: fixed edges on both sides would place
this vertex between two already exposed vertices.  There are one or
two such edges.  If no unexposed successor is available, choose
dually a maximal unexposed predecessor.  It has one or two fixed
outgoing edges.  Continue until the components touching the exposed
set are exhausted.  Adjacency guarantees that one of these choices is
available whenever such a component remains.  The local formulas for
both permitted singleton types apply, including to a nonscattering
contact.  Given top ends can be regarded as fictitious exposed
vertices later than all contacts on their lines.
For this argument one may use literal cuts in the uncontracted
four-edge coarea representation
\eqref{eq:atom-C}--\eqref{eq:edge-integral}.
Tonelli's theorem identifies each split bond variable with the
given boundary value of the next piece.  The activity exponents
add because a split bond is replaced by one integrated free end
and one fixed end.  The final singleton pieces are regular:
their fixed ends are on the same time side, so no two of them
are serial.  We do not assert that every intermediate complement
of these literal cuts is regular under the serial-segment
cutting convention described above.
\end{proof}

\subsection{A strengthened estimate for adjacent scattering contacts}

Write \(R\geq2\) for a polylogarithmic bound on velocities and on the
radii of the localized position sets.  The centers of the position
sets may be arbitrary.  Powers of \(R\) below may change from line to
line.  The time interval has length at most \(T\).
A pair of adjacent contacts with one bond is called \emph{primitive}
if there is no other directed chain between them.
This condition makes the pair order-convex.
A double bond between two consecutive scattering contacts would
require the same pair of freely separating spheres to collide again
without an intervening velocity change.  Such a history is impossible.
We therefore need only single-bond pairs in a realizable graph.

\begin{lemma}[Local adjacent-pair estimates]\label{lem:cutoff-pair}
Consider every possible choice $e_1,e_2$ of the two prospective fixed
edges of a type-A pair, with $t_1,t_2$ the corresponding contact times.
Call that choice strongly close if
\begin{equation}\label{eq:cutoff-close}
 |t_2-t_1|\leq\eps^{9/10},\qquad
 |v_{e_2}-v_{e_1}|\leq\eps^{4/5}.
\end{equation}
Then:
\begin{enumerate}[label=\textup{(\roman*)}]
\item a primitive pair satisfying \eqref{eq:cutoff-close}, cut
freely as one piece, has mass at most
\begin{equation}\label{eq:cutoff-free-pair}
 C R^C\mu\eps^{9/10+3(4/5)}
 =C R^C\eps^{13/10};
\end{equation}
\item in the absence of every strongly close primitive pair, each
type-A pair used by a cutting sequence has conditional mass at most
\(C R^C\eps^{1/20}\), or may be split into a neutral type-two
piece and a singleton with this same improved bound.
\end{enumerate}
The conditional bounds are uniform in the previously exposed lines
and times.
\end{lemma}

\begin{proof}
Reverse time within the local calculation if necessary.  In the
type-A arrangement the second contact has both outgoing edges
available, and its fixed incoming edge is not the shared bond.
Integrate the first contact as a type-three singleton and the second
as a type-two singleton.  The free first-contact parameters are
\((t_1,\omega,w)\).  Denote the two given velocities by \(u,v\),
and the intervening bond velocity by \(z\).

For the free pair in (i), use the two-contact tree parameters:
one free collision position, the two velocities at the first
contact, the new incident velocity at the second contact, two
normals, and two times.  The activity prefactor is \(\mu\).
Elastic reversal preserves velocity Lebesgue measure and the flux,
so this parametrization also covers a prospective fixed edge
chosen on the outgoing side of the first contact.  Conditional on
the first contact, the second given velocity is the independent
incident velocity of the third particle.  Its permitted set in
\eqref{eq:cutoff-close} has volume \(C\eps^{12/5}\); the second
time has interval length at most \(2\eps^{9/10}\).
All remaining localized volumes and fluxes cost at most \(CR^C\).
This proves \eqref{eq:cutoff-free-pair}.  By primitivity,
\Cref{lem:cutoff-neutral} finishes the rest of the connected graph
without another free singleton.

For (ii), first suppose \(t_2-t_1>\eps^{9/10}\).
Let \(q_1(t)\), \(q_2(t)\) be the positions on the given affine
lines.  The two contact equations imply
\begin{equation}\label{eq:cutoff-aiming}
 (t_2-t_1)(z-v)=q_2(t_1)-q_1(t_1)+O(\eps).
\end{equation}
The \(O(\eps)\) error includes both contact offsets and is uniform
in the normals.  At fixed \(t_1\), the vector \(z-v\) consequently
lies within distance \(C\eps^{1/10}\) of the line
\(\R(q_2(t_1)-q_1(t_1))\).  If this displacement is zero it instead
lies in a ball of that radius, which is a stronger restriction.

Put \(h=C\eps^{1/10}\).  Depending on which edge is fixed at the
first contact, the bond velocity is one of
\[
 w,\qquad
 P_{\omega^\perp}w+(u\cdot\omega)\omega,\qquad
 u+s\omega,\quad s=(w-u)\cdot\omega ,
\]
up to signs that do not change these estimates.
For the first map, the preimage of an \(h\)-tube has
three-dimensional volume at most \(CR^C h^2\).
For the second, a plane intersects a bounded part of the tube
in area at most \(CR^C h\), and the unused longitudinal
coordinate costs only \(CR\).
For the third, integrating the two tangential coordinates of \(w\)
costs \(CR^2\), while the image of \(\dd s\,\dd\omega\) has
density at most \(C|z-u|^{-2}\).
Inside \(|z-u|\leq\sqrt h\) its mass is at most \(C\sqrt h\).
Outside that ball, the density is at most \(Ch^{-1}\) and the
tube volume is at most \(CRh^2\).
Thus its total mass is at most \(CR^C\sqrt h\).
These bounds, with the first flux bounded by \(CR\) and the
type-two kernel by one, give \(CR^C\eps^{1/20}\).

Next suppose \(t_2-t_1\leq\eps^{9/10}\) and
\(|u-v|>\eps^{4/5}\).
Equation \eqref{eq:cutoff-aiming} gives
\[
 |q_2(t_1)-q_1(t_1)|\leq CR\eps^{9/10}.
\]
The difference of the given affine lines has velocity \(v-u\);
the set of such \(t_1\) is an interval of length at most
\(CR\eps^{1/10}\).  Integrating this restricted first time,
and bounding all other local factors by \(CR^C\), gives a
stronger estimate.

Finally suppose \eqref{eq:cutoff-close} holds.  The pair is then
nonprimitive.  Choose another directed chain between its vertices.
At least one intermediate vertex of this chain has already been
cut.  Indeed, if all were still uncut, the first vertex would have
two uncut outgoing bonds and the second two uncut incoming bonds.
Their fixed edges would be on the opposite outer sides, which is
precisely not the type-A arrangement.  A time \(t_*\) strictly
between \(t_1,t_2\) is therefore part of the previously exposed
data.  The first free collision time lies within
\(\eps^{9/10}\) of \(t_*\).  Cutting the pair into a type-three
and a type-two singleton gives the claimed estimate.
There are only polynomially many choices of such intermediate
vertices in a fixed graph; a deterministic choice suffices.
\end{proof}

In the sequel fix the dimensional number
\begin{equation}\label{eq:cutoff-delta}
 \delta=\tfrac12\min\{\upsilon,1/20\}>0.
\end{equation}
Every retained local gain, whether from
\Cref{lem:cutoff-pair} or \Cref{lem:dhm-local}, is therefore
at least a factor \(\eps^\delta\), apart from a fixed
polylogarithmic factor.

\subsection{Accumulating a fixed amount of gain}

\begin{lemma}[Amplified pure-collision cut]\label{lem:cutoff-amplified}
For a sufficiently large dimensional \(\Gamma_0\), every connected
realizable full scattering graph with
\(\Gamma<r(M)\leq2\Gamma\), \(\Gamma\geq\Gamma_0\), has a
finite decomposition into cut sequences of the following form.
Either it contains a freely cut pair estimated by
\eqref{eq:cutoff-free-pair} and no other activity loss, or its
retained improved pieces and free singletons satisfy
\begin{equation}\label{eq:cutoff-gain-balance}
 \delta\,g-2s_4\geq6/5 .
\end{equation}
Here \(g\) counts the retained improved pieces, and \(s_4\)
counts free singleton pieces.  The number of subcases is at most
\(C^b|\log\eps|^{C_*}\).  The number of pieces other than
neutral type-three singletons is bounded in terms of \(\Gamma\)
and the fixed choices.
\end{lemma}

\begin{proof}
First take the union over all strongly close primitive pairs and
all prospective type-A fixed-edge choices.  There are at most
a constant times \(b\) candidate bonds and a bounded number of
edge choices per bond.  A fixed priority partitions this union.
Cut the selected pair freely and use
\Cref{lem:cutoff-neutral}.  This is the first alternative.
On the complement every type-A pair has the gain
in \Cref{lem:cutoff-pair}.

Here is an explicit finite hierarchy for the second alternative.
Choose
\[
 D=\lceil4/\delta\rceil,\qquad
 B=100/\delta,\qquad
 R_0=(10^5+100/\delta)(1+2B),\qquad
 \kappa_1=2,
\]
and choose increasing integers recursively so that
\[
 \kappa_{j+1}>
 \kappa_j+40\{1+G(\lceil R_0\kappa_j\rceil)\},
 \qquad 1\leq j<D .
\]
Use an increasing envelope of \(G\) if needed at its two
smallest arguments.  Take \(\Gamma_0\geq\kappa_D\).
All these choices precede \(L\).

Apply the transversal chain of \Cref{prop:cutoff-import}.
If it has at least \(D\) transitions, its transition sequence
has \(g\geq D\), \(s_4=1\), and
\(\delta D-2\geq2>6/5\).
Otherwise the chain reaches \(M\) in at most \(D\) sets.
Since \(r(A_1)=0<\kappa_1\) and
\(r(M)>\Gamma\geq\kappa_D\), there is a first \(j\) such that
\[
 r(A_j)<\kappa_j,\qquad r(A_{j+1})\geq\kappa_{j+1}.
\]
Put \(A=A_j\).  The growth bound shows that one of
\(X^\pm(A)\) has size greater than
\(G(\lceil R_0\kappa_j\rceil)\).
Its number \(m\) of connections to \(A\) therefore satisfies
\(m>R_0\kappa_j\), by
\Cref{prop:cutoff-import}\ref{item:cutoff-bfk}.
The retained realizability of the witness is used at exactly
this invocation of the collision-count bound.

For the chosen side put \(z=|X_0^\pm(A)|\), \(a=r(A)\),
and use the partition and integer \(y\) of
\Cref{prop:cutoff-import}\ref{item:cutoff-main}.
If \(z\geq B\kappa_j\), use the transversal sequence:
\[
 \delta g-2s_4
 \geq \delta\{(B/2-1)\kappa_j-1\}-2>6/5 .
\]
If \(y\geq B\kappa_j\), use its first main sequence:
\[
 \delta g-2s_4
 \geq \delta\{(B/10-3)\kappa_j-3\}-2>6/5 .
\]
Otherwise \(S=a+z+y\leq(1+2B)\kappa_j\).
In the second main sequence, every counted type-A or type-B
pair has a gain.  Hence
\[
 \delta g-2s_4
 \geq \frac{\delta}{10}(m-10^5S)-2S
 > 8(1+2B)\kappa_j>6/5 .
\]
If more uniformly good singletons are produced than needed,
retain only the number of their gains needed for this inequality.

For completeness, a degree count bounds the remaining exceptional
pieces.  Write \(s_j\) for singleton counts and \(p_{33},p_{44}\)
for pair counts.  Counting each bond once, except that an internal
pair bond contributes at both endpoints when available degrees
are summed, gives
\begin{equation}\label{eq:cutoff-degree-count}
 s_2+p_{33}=r(M)+s_4+p_{44}-1 .
\end{equation}
One way to verify this identity is to use
\(b=s_2+s_3+s_4+2p_{33}+2p_{44}\) and
\(n=b+1-r(M)\).  The sum of available degrees is both
\[
 2s_2+3s_3+4s_4+6p_{33}+8p_{44}
 \quad\text{and}\quad
 2b+n+p_{33}+p_{44}.
\]
This proves \eqref{eq:cutoff-degree-count}.
The constructed \(s_4\) is bounded, \(r(M)\leq2\Gamma\),
and \(p_{44}\) is zero except for the single primitive seed.
Thus \(s_2,p_{33}\) are bounded.  Splitting a nonprimitive
type-A pair as in \Cref{lem:cutoff-pair} changes these counts
by a bounded amount.  The number of retained good singletons
may also be bounded in terms of the same constants.
The subcase count follows from
\Cref{prop:cutoff-import} and the additional polynomial choices.
\end{proof}

\subsection{Simultaneous integration of the ordinary parameters}

The local powers must be retained without a logarithmic loss at each
ordinary contact.  We give the parameter argument in detail.  In a
cut sequence call every neutral type-three singleton \emph{ordinary};
all other pieces, including the retained good singletons, are
\emph{exceptional}.  Counting both vertices of an exceptional pair,
their number is bounded by a constant depending on \(K,\Gamma\).
For the terminal part this follows from
\Cref{lem:cutoff-amplified}.  Below it use
\Cref{lem:cutoff-neutral}.  The identity
\eqref{eq:cutoff-degree-count} on the complete graph, or its version
with the bounded number of crossing ends fixed, bounds the additional
type-two pieces.  Empty fixed lines require only deterministic
substitution; any such lines can be included in the exceptional list.
There is no full empty line, since every line is attached to the
terminal component.

Let \(q\) denote the number of ordinary atoms and \(b\) the total
number of atoms.  Thus
\begin{equation}\label{eq:cutoff-ordinary-count}
 q=b-O_K(1).
\end{equation}
The constants in exponential factors below will not depend on \(K\).
Dependence on \(K\) is allowed in prefactors and in fixed polynomial
degrees.

We will bound the simultaneous velocity domain using total birth
energy, and trade the causal time restrictions against the ordinary
relative-speed factors. Both bounds concern the original feasible
histories, before any remaining collision constraints are dropped.

\begin{lemma}[Extraction of literal free parameters]\label{lem:cutoff-extraction}
Consider a finite ordered composition of positive cut kernels.
At an ordinary step its free variables are
\(\vartheta_j=(t_j,\omega_j,w_j)\), with measure
\[
 a_j\,\dd t_j\,\dd\omega_j\,\dd w_j ,
\]
followed by a deterministic update of the exposed data.
Suppose every exceptional kernel \(K_\ell(h,\dd z)\) has mass
at most \(H_\ell\), uniformly in admissible data \(h\) from
earlier cuts.  Include in that kernel every restriction needed
for this bound, and set it to zero outside its admissibility
domain.

Remove the factors \(a_j\) from the ordinary measures, obtaining
a positive measure \(\nu\).  If \(E\) is a set of complete
histories and \(\mathcal D\) is a measurable set containing its
projection onto the literal variables
\(\vartheta=(\vartheta_1,\ldots,\vartheta_q)\), then
\begin{equation}\label{eq:cutoff-extraction}
 \nu(E)\leq
 \left(\prod_\ell H_\ell\right)
 \int_{\mathcal D}\prod_{j=1}^q
       \dd t_j\,\dd\omega_j\,\dd w_j .
\end{equation}
The same assertion holds when the ordinary density after division
by \(a_j\) is merely bounded by one.
\end{lemma}

\begin{proof}
First put all ordinary variables in bounded boxes.  Fix their
entire vector \(\vartheta\).  Their deterministic updates still
occur in their original positions in the cutting sequence.
The exceptional kernels remain in their original order, with
\(\vartheta\) treated as exterior parameters.  Integrating these
kernels backwards from the last exceptional step bounds the
resulting total mass by \(\prod_\ell H_\ell\).
Intervening deterministic substitutions have norm at most one
on bounded functions.  Nonnegative kernel Fubini therefore
represents \(\nu\) as the product measure in \(\vartheta\)
times these remaining conditional kernels.  Majorizing the
indicator of \(E\) by \(\ind_{\mathcal D}(\vartheta)\) proves
\eqref{eq:cutoff-extraction}.  Exhaustion removes the bounded
boxes.

In particular, this calculation does not condition an exceptional
kernel on the global energy event.  That event is retained until
its ordinary projection is majorized.  Fixing a future ordinary
parameter does not alter any exceptional kernel or require an
averaged estimate to become pointwise.  Uniformity in the actual
earlier-cut data is sufficient.

There is also a rectangle formulation useful for the coarea
measures.  Restrict the ordinary relative speeds away from zero
and bound their product on the original feasible support.
After division by those speeds, an ordinary local kernel in a
rectangle of its \((t,\omega,w)\) coordinates has mass bounded
by the rectangle's product volume.  In a product rectangle
for all ordinary steps, successive conditional estimates give
that volume times \(\prod_\ell H_\ell\).  Cover a compact
containing set \(\mathcal D\) by finite unions of product
rectangles with volume decreasing to its outer measure,
using finitely many sphere charts.  Summing the rectangle
bounds and taking the infimum gives
\eqref{eq:cutoff-extraction}.  The containing sets used below
are compact and measurable.  Finally let the speed restriction
decrease to zero; a zero-speed incoming contact has zero flux.
This proves the same statement directly by exhaustion of
transversal coarea histories.
\end{proof}

The literal extraction in this lemma is the operator rearrangement
in \cite[Equations (9.66), (9.74)--(9.77)]{DHM2025}.
Its use with the following containing sets, rather than a
separate velocity box at every ordinary vertex, is the additional
estimate needed here.

\begin{lemma}[A common velocity containing set]\label{lem:cutoff-ellipsoid}
Fix one original history graph and its ordinary normals.
Suppose the sum of squared velocities at all its births is at
most \(P\), and suppose there are \(s\) exceptional atoms,
counting both atoms of a pair.  For \(q\geq1\), the concatenated
ordinary parameters \(W=(w_1,\ldots,w_q)\) lie in a set of
volume at most
\begin{equation}\label{eq:cutoff-ellipsoid}
 C_s\{(1+P)(1+q)\}^{C s}
 \left(C\sqrt{1+P/q}\right)^{3q}.
\end{equation}
The exponential base \(C\) is absolute.  The containing set
depends on the ordinary normals and the graph but not on the
numerical contact times or on the exceptional parameters.
If \(P,q\) are bounded by fixed powers of \(|\log\eps|\),
the prefactor is a fixed power of \(|\log\eps|\).
\end{lemma}

\begin{proof}
Work on the original feasible histories, before dropping any
collision constraints.  Let \(u\) be the vector of all birth
velocities.  At each exceptional atom replace its two outgoing
velocities by independent additional inputs, and write \(a\)
for their concatenation.  There are at most \(6s\) scalar
coordinates in \(a\).  On the energy support each of its
three-dimensional entries has norm at most \(\sqrt P\).

Run the resulting velocity computation forward in the graph's
chronological order.  Ordinary scattering applies an orthogonal
matrix to two particle velocities; an ordinary nonscattering
contact applies the identity.  At an exceptional atom, erase
the dependence of its two outgoing coordinates on \(u\), and
insert the prescribed values from \(a\).  Birth coordinates
can be present from the start as inactive coordinates, and
departed coordinates can be retained as inactive ones.
The part depending linearly on \(u\) is thus propagated by
orthogonal matrices and coordinate projections.  Its operator
norm is at most one at every time.

Every ordinary parameter \(w_j\) is one of the physical edge
velocities immediately before or after its contact.  This
includes a parameter on the outgoing side when the cut
parametrization runs backwards.  Consequently
\begin{equation}\label{eq:cutoff-forward-rows}
 W=A(\omega)u+d(\omega,a),\qquad
 \|A(\omega)\|_{\mathrm{HS}}^2\leq3q .
\end{equation}
Repeated observation of one velocity does not require
orthogonality between different rows; the stated sum of
their squared norms is all that is used.
These maps do not depend on numerical contact times.
Different total extensions of the causal order produce the
same maps, since incomparable contacts act on disjoint
particle coordinates and commute.

Take a Euclidean net of the possible \(a\)'s of mesh
\(c_s(1+q)^{-1/2}\).  It has cardinality at most
\(C_s\{(1+P)(1+q)\}^{Cs}\).
A change of size \(\zeta\) in this entire vector changes
any one chronological velocity state by at most
\(\sqrt s\,\zeta\): propagate the reset errors and use
Cauchy--Schwarz.  Its effect on the \(q\) observations is
at most \(C\sqrt{qs}\,\zeta\).
Choose \(c_s\) so that the simultaneous error is at most one.
The allowed \(W\)'s are contained in the union, over net
points \(a_0\), of
\begin{equation}\label{eq:cutoff-ellipsoid-set}
 d(\omega,a_0)+A(\omega)B_{3n}(\sqrt P)+B_{3q}(1).
\end{equation}

For \(m=3q\) the Minkowski sum in
\eqref{eq:cutoff-ellipsoid-set} is contained in
\[
 \sqrt2\,(PAA^{\mathsf T}+\Id_m)^{1/2}B_m(1).
\]
Indeed its support function in direction \(\xi\) satisfies
\[
 \sqrt P|A^{\mathsf T}\xi|+|\xi|
 \leq \sqrt{2(P|A^{\mathsf T}\xi|^2+|\xi|^2)}.
\]
Combine the trace bound in \eqref{eq:cutoff-forward-rows}
with the arithmetic--geometric mean inequality for eigenvalues
and the ball-volume estimate
\[
 \operatorname{vol}(B_m(1))\leq(C/\sqrt m)^m.
\]
This gives
\[
 2^{m/2}\operatorname{vol}(B_m(1))
 \det(\Id_m+PAA^{\mathsf T})^{1/2}
 \leq \left(C\sqrt{(1+P)/q}\right)^{3q}
 \leq \left(C\sqrt{1+P/q}\right)^{3q}.
 \]
Sum over the net.  The identity-matrix enlargement handles
zero or small singular values without any inverse Jacobian.
For varying normals, the sets in
\eqref{eq:cutoff-ellipsoid-set} form a finite union of
continuous images of compact parameter sets.  They provide
a compact measurable containing set for use in
\Cref{lem:cutoff-extraction}.
\end{proof}

\begin{lemma}[Time order and relative speeds]\label{lem:cutoff-time}
In a slab containing \(q_i\) ordinary atoms, retain the partial
order induced on those atoms by the full original graph,
including chains through exceptional atoms.  Let
\(\mathcal T_i\) be its time-order region in that slab.
On the original support with total birth energy at most \(P\),
\begin{equation}\label{eq:cutoff-time}
 |\mathcal T_i|\,
 \sup \prod_{a\text{ ordinary in slab }i}|v_a-v_{*a}|
 \leq
 \frac{\tau^{q_i}(2q_iP)^{q_i/2}}{q_i!}.
\end{equation}
An empty slab contributes the factor one.
\end{lemma}

\begin{proof}
Let \(e_i\) be the number of linear extensions of this partial
order.  Apart from equality hyperplanes, its time-order
region is the disjoint union of \(e_i\) simplexes, so
\(|\mathcal T_i|=e_i\tau^{q_i}/q_i!\).

An antichain of ordinary vertices occupies disjoint incoming
pairs on a common causal slice.  To see this, place all strict
ancestors of the antichain first in a total extension of the
full graph.  No two vertices in the antichain share a particle
line.  Their incoming states are present simultaneously on
the resulting slice.  Conservation at all original contacts,
including exceptional contacts, bounds the sum of their pair
energies by the total birth energy \(P\).  Retaining unborn
and departed velocities as inactive coordinates makes the
same argument valid in the presence of interfaces.

For a fixed feasible velocity assignment write
\(a_v=|v_v-v_{*v}|\).  Every antichain \(\mathcal A\)
therefore satisfies
\[
 \sum_{v\in\mathcal A}a_v
 \leq\sqrt{2|\mathcal A|P}\leq\sqrt{2q_iP}.
\]
In the recursive enumeration of linear extensions, the
available minimal vertices are an antichain at each step.
Summing the products of their weights successively gives
\[
 e_i\prod_v a_v\leq(2q_iP)^{q_i/2}.
\]
This bound holds for every feasible assignment, so also
for the supremum of its product.  Multiplication by
\(\tau^{q_i}/q_i!\) proves \eqref{eq:cutoff-time}.
In particular, no assumption that the maximizing assignment
is independent of the times is necessary.
\end{proof}

\begin{lemma}[Summation with finitely many exceptional pieces]\label{lem:cutoff-sum}
Consider the positive low-complexity graph integrals obtained
above.  Suppose a cut sequence has \(q=b-O_K(1)\) ordinary
atoms and boundedly many exceptional pieces, whose conditional
mass bounds have product
\(\eps^\alpha|\log\eps|^{C_*}\).
After summing spatial cells and energy bins in a polylogarithmic
energy region, its integral is at most
\begin{equation}\label{eq:cutoff-sum}
 \eps^\alpha|\log\eps|^{C_*}
       (C\tau\sqrt L)^q .
\end{equation}
The exponential base \(C\) depends on the fixed kinetic
envelopes and dimension, but not on \(K,L\).
\end{lemma}

\begin{proof}
By \Cref{prop:one-root,lem:kinetic-envelopes}, choose
\(\beta_*>0\) and a common Gaussian Bol bound for all birth
inputs, with \(\beta_*\) independent of \(L\).  For each
birth factor use its spatial-cell decomposition
\[
 |g_i(x,v)|\leq C\sum_{a\in\Z^3}m_{i,a}
       \ind_{\{|x-a|\leq C\}}e^{-\beta_*|v|^2},
 \qquad \sum_a m_{i,a}\leq C .
\]
The same form applies at time zero.  Sum over the birth-cell
indices at the end; their coefficients have total at most
\(C^n=C^{b+O_K(1)}\).

Decompose total birth energy into \([0,1]\) and the dyadic
bins \([P/2,P]\), \(P\geq1\).  On a bin, the Gaussian factors
give \(C^{b+O_K(1)}e^{-cP}\).  Fixed birth cells and bounded
energy localize each physical particle's positions within
a polynomially bounded distance of its birth cell over
\([0,T]\).  The same is true of its affine intercept.
Thus each exceptional kernel has the position localization
required by its stated uniform estimate, with an arbitrary
but fixed cell-dependent center.

For the ordinary steps use the relative speed rather than
the smaller normal flux as \(a_j\) in
\Cref{lem:cutoff-extraction}.  Its quotient density is at
most one.  Pull out the product of the per-slab suprema
in \Cref{lem:cutoff-time}, on the original feasible support.
Use the containing velocity sets of
\Cref{lem:cutoff-ellipsoid}, and keep the full ordinary
time orders.  All gain-producing exceptional restrictions
remain attached to their kernels.  Other future restrictions
may now be dropped when taking their uniform mass bounds.
The ordinary sphere integration contributes \((4\pi)^q\).
After absorbing the bounded exceptional-net prefactor into
the logarithmic power, the bin is bounded by
\[
 \eps^\alpha|\log\eps|^{C_*}
 C^q e^{-cP}(1+P/q)^{3q/2}
 \prod_i\frac{\tau^{q_i}(2q_iP)^{q_i/2}}{q_i!}.
\]
When \(q\geq1\), Stirling's inequality and convexity give
\[
 q_i!\geq(q_i/e)^{q_i},\qquad
 \prod_iq_i^{q_i}\geq(q/L)^q ,
\]
with zero counts omitted.  The last bound becomes
\begin{equation}\label{eq:cutoff-energy-bin}
 \eps^\alpha|\log\eps|^{C_*}
 (C\tau\sqrt L)^q e^{-cP}
 (P/q)^{q/2}(1+P/q)^{3q/2}.
\end{equation}
For \(x=P/q\),
\[
 \sup_{x\geq0}x^{1/2}(1+x)^{3/2}e^{-cx/2}<\infty .
\]
Consequently the final three factors in
\eqref{eq:cutoff-energy-bin} are bounded by
\(C^qe^{-cP/2}\).  Summing the dyadic bins costs only
another absolute constant.  The lowest bin is treated
by the same bound with \(P=1\).
If \(q=0\), integrate the bounded number of exceptional
pieces directly; no time factor is asserted for them.
This proves \eqref{eq:cutoff-sum}.
\end{proof}

\subsection{Summing the terminal witnesses}

\begin{proof}[Proof of \Cref{prop:cutoff}]
Use the terminal witness reduction and choose \(K\) as in
\eqref{eq:cutoff-low-complexity}.  Apply
\Cref{lem:cutoff-energy-tail}.  For the remaining graphs all
velocities are polylogarithmically bounded.  Decomposing the
birth positions into cells as in
\Cref{lem:cutoff-sum} also gives the required local position
regions: along each original particle line the physical
travel distance is at most \(T\sqrt{P_\eps}\); its affine
intercepts are within another \(T\sqrt{P_\eps}\).
Their centers are fixed by the birth cells and may be
arbitrary in the local estimates.

First consider witnesses with
\(\Gamma<r(M_{\mathrm{top}})\leq2\Gamma\).
Cut the terminal component freely and apply
\Cref{lem:cutoff-amplified}. The terminal atom set is order-convex:
every lower atom lies in an earlier slab, while a directed path
has increasing times. Straight continuations beyond the stopping
contact add no atoms. Finish all lower components
from their now fixed terminal data by
\Cref{lem:cutoff-neutral}.  These lower components create
no extra type-four singleton: by \Cref{cor:terminal-witness},
each is attached to the terminal part through particle lines,
without using initial links. The total cycle and crossing bounds
give \eqref{eq:cutoff-ordinary-count} for the entire cut
sequence.

In the primitive branch the product of exceptional masses is
at most \(\eps^{13/10}|\log\eps|^{C_*}\).
In every other branch it is at most
\(\eps^{6/5}|\log\eps|^{C_*}\), by
\eqref{eq:cutoff-gain-balance}; a type-four singleton costs
\(\mu=\eps^{-2}\).  The positions, exceptional variables,
and intermediate-time choices only enlarge the fixed
logarithmic exponent.  Therefore
\Cref{lem:cutoff-sum} bounds either branch by
\[
 \eps^{6/5}|\log\eps|^{C_*}
       (C\tau\sqrt L)^{b-O_K(1)}.
\]

We explain explicitly the uniformity needed to sum these
bounds over graphs.  For bounded \(\rho\), the unlabelled
root and interface counting in
\Cref{thm:dhm-package,cor:terminal-witness}
costs \(C_1^b\) times a fixed power of the cutoffs.
The pure-collision cutting partitions have the same form
by \Cref{prop:cutoff-import}.  The additional selected
pairs, crossing labels, and initial coefficients have
boundedly many choices of locations, and therefore cost
only fixed-degree polynomials in the graph size.
The exponential bases depend on \(\Gamma\) and the
uniform kinetic envelopes; they do not depend on \(K\)
or \(L\).  The \(K\)-dependence is confined to the fixed
polynomial and logarithmic exponents.
To see explicitly why assigning slabs does not add a factor
\(L^b\), cut the bounded number of interface bonds and delete
the bounded number of cycle bonds.  The remaining atom forests
have degree at most four.  Ordered-tree encodings and the finite
local C/O and serial-edge choices cost \(C^b\); restoring bonds
and marked ends costs \((1+b)^{O(K+\Gamma)}\).
Each resulting component lies in one slab, so its slab choices
cost only \(L^{O(K+\Gamma+1)}\), a front constant.
Equivalently, first choose the boundedly many occupied slab
indices, then count the positive size compositions by \(2^b\).
Empty continued lines are bounded by the root and crossing counts.
Choose \(L\) so large that
\[
 C_1 C\tau\sqrt L=C_1CT/\sqrt L<1/2 .
\]
The \(O_K(1)\) missing ordinary factors contribute a
constant depending on \(K,L\), and the geometric series
in \(b\) converges.  The sum of this first class is
\begin{equation}\label{eq:cutoff-cycle-sum}
 O\bigl(\eps^{6/5}|\log\eps|^{C_*}\bigr).
\end{equation}

For the other witnesses use
\(n(M_{\mathrm{top}})>\Lambda_i\) and
\(r(M_{\mathrm{top}})\leq2\Gamma\).
Choose one terminal vertex as a free type-four seed and
finish by \Cref{lem:cutoff-neutral}; complete the lower
part in the same manner.  There is one activity loss
\(\mu\), and no gain is required.
Degree counting again leaves \(b-O_K(1)\) ordinary
vertices.  Moreover
\(b\geq b(M_{\mathrm{top}})
       =n(M_{\mathrm{top}})-1+r(M_{\mathrm{top}})
       \geq\Lambda_i-1\).
The same summation now gives
\[
 C_*\eps^{-2}|\log\eps|^{C_*}
       \sum_{b\geq\Lambda_i-1}2^{-b},
\]
which is smaller than every fixed power of \(\eps\),
since \(\Lambda_i\geq|\log\eps|^{30}\).

Finally sum over the finitely many slabs and add the
coarse high-complexity, high-energy, and expansion
remainders.  For sufficiently small \(\eps\),
\eqref{eq:cutoff-cycle-sum} is bounded by
\(C\eps^{1+1/20}\).  This proves
\eqref{eq:cutoff-probability}.

The order of choices is worth recording.  The local exponent
\(\delta\) and the finite hierarchy in
\Cref{lem:cutoff-amplified} fix \(\Gamma_0\) using only
dimensional constants.  Choose \(\Gamma\), then choose
\(L\) for the imported estimates and the common exponential
base.  Only afterwards choose the complexity threshold
\(K\), localization powers, and the sufficiently small
upper bound on \(\eps\).  No observation number or moment
order enters the first two choices.
\end{proof}
